\documentclass[12pt]{article}
\usepackage{amsmath, amssymb}
\usepackage[margin=1in]{geometry}
\setlength{\parskip}{6pt}
\title{QUBO Formulation: Vehicle Routing Assignment Problem}
\date{}
\begin{document}
\maketitle

\subsection*{Vehicle Routing Assignment Problem}

\paragraph{Problem Statement.}
Given a set of customers $\mathcal{I}$, a set of vehicles $\mathcal{K}$, and a set of candidate routes $\mathcal{R} = \{0, 1, \dots, M-1\}$, where each route $j \in \mathcal{R}$ has an associated non-negative cost $c_j$ and serves a subset of customers $C_j \subseteq \mathcal{I}$ while requiring a subset of vehicles $V_j \subseteq \mathcal{K}$, the task is to select a subset of routes that minimizes the total route cost subject to the constraints that every customer is served exactly once and each vehicle is assigned to at most one route.

\paragraph{Decision Variables.}
For each candidate route $j \in \mathcal{R}$, we introduce a binary decision variable $x_j \in \{0, 1\}$. The variable $x_j$ takes the value $1$ if route $j$ is selected, and $0$ otherwise. The complete solution is represented by the vector $x = (x_0, x_1, \dots, x_{M-1})^\top \in \{0, 1\}^M$.

\paragraph{QUBO Derivation.}
We construct a Quadratic Unconstrained Binary Optimization (QUBO) Hamiltonian $H(x)$ by converting the constrained optimization problem into an unconstrained energy minimization problem using the penalty method.

\textbf{Step 1 --- Objective Function.}
The original objective is to minimize the total cost of selected routes. Since the QUBO framework naturally performs minimization, we retain the original sign:
\begin{align}
  H_{\text{obj}}(x) = \sum_{j \in \mathcal{R}} c_j x_j.
\end{align}

\textbf{Step 2 --- Constraints.}
The problem is subject to two types of constraints:
\begin{align}
  \sum_{j \in \mathcal{R}: i \in C_j} x_j = 1, \quad \forall i \in \mathcal{I}, \\
  \sum_{j \in \mathcal{R}: k \in V_j} x_j \le 1, \quad \forall k \in \mathcal{K}.
\end{align}

\textbf{Step 3 --- Penalty Terms.}
We enforce the constraints by adding squared penalty terms to the objective. Let $\lambda_c > 0$ and $\lambda_v > 0$ be penalty weights.
For the equality constraint, the penalty for customer $i$ is:
\begin{align}
  P_i^{\text{cust}}(x) = \lambda_c \left( \sum_{j \in \mathcal{R}: i \in C_j} x_j - 1 \right)^2.
\end{align}
Expanding the square and applying the idempotent property $x_j^2 = x_j$ for binary variables yields:
\begin{align}
  P_i^{\text{cust}}(x) &= \lambda_c \left( \sum_{j \in \mathcal{R}: i \in C_j} x_j^2 - 2 \sum_{j \in \mathcal{R}: i \in C_j} x_j + 1 + \sum_{\substack{j, l \in \mathcal{R}: i \in C_j \\ j \neq l}} x_j x_l \right) \nonumber \\
  &= \lambda_c \left( 1 - \sum_{j \in \mathcal{R}: i \in C_j} x_j + 2 \sum_{\substack{j < l \\ j, l \in \mathcal{R}: i \in C_j}} x_j x_l \right).
\end{align}
For the inequality constraint, since $x_j \in \{0, 1\}$, the condition $\sum x_j \le 1$ is violated if and only if at least two routes assigned to vehicle $k$ are simultaneously selected. We penalize every violating pair:
\begin{align}
  P_k^{\text{veh}}(x) = \lambda_v \sum_{\substack{j < l \\ j, l \in \mathcal{R}: k \in V_j}} x_j x_l.
\end{align}

\textbf{Step 4 --- Combined Hamiltonian.}
Summing the objective and all penalty contributions gives the total QUBO Hamiltonian:
\begin{align}
  H(x) = \sum_{j \in \mathcal{R}} c_j x_j + \sum_{i \in \mathcal{I}} \lambda_c \left( 1 - \sum_{j \in \mathcal{R}: i \in C_j} x_j + 2 \sum_{\substack{j < l \\ j, l \in \mathcal{R}: i \in C_j}} x_j x_l \right) + \sum_{k \in \mathcal{K}} \lambda_v \sum_{\substack{j < l \\ j, l \in \mathcal{R}: k \in V_j}} x_j x_l.
\end{align}

\textbf{Step 5 --- Q Matrix and Offset.}
Collecting coefficients for linear and quadratic terms, we identify the diagonal entries $Q_{jj}$, off-diagonal entries $Q_{jl}$ (for $j < l$), and the constant offset $c$:
\begin{align}
  Q_{jj} &= c_j - \lambda_c \sum_{i \in \mathcal{I}} \mathbb{I}(i \in C_j), \quad \forall j \in \mathcal{R}, \\
  Q_{jl} &= \lambda_c \sum_{i \in \mathcal{I}} \mathbb{I}(i \in C_j \cap C_l) + \lambda_v \sum_{k \in \mathcal{K}} \mathbb{I}(k \in V_j \cap V_l), \quad \forall j < l, \\
  c &= \lambda_c |\mathcal{I}|,
\end{align}
where $\mathbb{I}(\cdot)$ denotes the indicator function. The Hamiltonian is thus expressed in standard QUBO form as $H(x) = \sum_{j \in \mathcal{R}} Q_{jj} x_j + \sum_{j < l} 2 Q_{jl} x_j x_l + c$.

\paragraph{Penalty Weight Justification.}
The penalty weights are chosen as $\lambda_c, \lambda_v > \max_{j \in \mathcal{R}} c_j$. For the equality constraint, the maximum objective gain from violating $\sum x_j = 1$ is bounded by the largest route cost. For the pairwise inequality constraint $x_j + x_l \le 1$, the penalty must strictly exceed the maximum single-variable objective coefficient to ensure that any pair of simultaneously selected routes incurs an energy increase larger than any possible cost reduction. Setting $\lambda_c, \lambda_v > \max_j c_j$ satisfies both conditions, guaranteeing that infeasible configurations are energetically unfavorable.

\paragraph{Correctness Argument.}
Minimizing $H(x)$ over $\{0, 1\}^M$ recovers the optimal solution to the original constrained problem. First, any feasible assignment satisfies all constraints, causing all penalty terms to vanish and reducing $H(x)$ exactly to the original objective function. Second, any infeasible assignment incurs a penalty strictly greater than the maximum possible improvement in the objective function from any feasible assignment. Consequently, the global minimum of $H(x)$ must correspond to a feasible solution, and among all feasible solutions, it must minimize the original cost.

\paragraph{Illustrative Example.}
Consider a concrete instance with $M = 3$ candidate routes, one customer ($\mathcal{I} = \{1\}$), and one vehicle ($\mathcal{K} = \{1\}$). Suppose all three routes serve the customer and require the vehicle, i.e., $C_0 = C_1 = C_2 = \{1\}$ and $V_0 = V_1 = V_2 = \{1\}$. The route costs are $c_0, c_1, c_2$.
Instantiating the general formulas yields the following QUBO parameters:
\begin{align}
  Q_{00} &= c_0 - \lambda_c, \quad Q_{11} = c_1 - \lambda_c, \quad Q_{22} = c_2 - \lambda_c, \\
  Q_{01} &= \lambda_c + \lambda_v, \quad Q_{02} = \lambda_c + \lambda_v, \quad Q_{12} = \lambda_c + \lambda_v, \\
  c &= \lambda_c.
\end{align}
The resulting Hamiltonian is:
\begin{align}
  H(x) &= (c_0 - \lambda_c)x_0 + (c_1 - \lambda_c)x_1 + (c_2 - \lambda_c)x_2 \nonumber \\
  &\quad + 2(\lambda_c + \lambda_v)(x_0 x_1 + x_0 x_2 + x_1 x_2) + \lambda_c.
\end{align}
Evaluating $H(x)$ over all $2^3 = 8$ binary assignments reveals that feasible solutions (exactly one $x_j = 1$) yield energies $c_0, c_1, c_2$ respectively. Infeasible solutions (zero or two or three selected routes) incur penalties that strictly exceed any feasible cost, given $\lambda_c, \lambda_v > \max_j c_j$. Thus, the global minimum corresponds to selecting the route with the smallest cost. For instance, if $c_0 < c_1 < c_2$, the optimal solution is $x^* = (1, 0, 0)^\top$ with minimum energy $H(x^*) = c_0$.
\end{document}